\thispagestyle{plain}
\chapter*{Resumen}

% Versión en español del resumen
\guideinfo{El objetivo de este proyecto es el diseño e implementación de un modelo de \acrfull{ia} capaz de automatizar la segmentación semántica de los usos del suelo en España a partir de
las imágenes del satélite Sentinel-2. El enfoque metodológico propuesto implicó la extracción de las imágenes satelitales, mediante Google Earth Engine, y la descarga de la base de datos
con los usos de suelo clasificados a nivel nacional ofrecida por \acrfull{corine} Land Cover España. Estos datos se analizaron con la librería Rasterio con el fin de comprobar que estuvieran en el
mismo \acrfull{crs}. Adicionalmente rasterizamos la información vectorial dentro del geopackage y la proyectamos sobre el sistema EPSG:25830 para alinearlo con
las imágenes del satélite con resolución de 10 metros por píxel y separando los macro parches en parches más manejables de 50x50.
\newline
Para la creación de la red neuronal, se propusieron dos modelos distintos: El primero una \acrfull{cnn} que desemboca en un \acrfull{mlp} con el fin de 
clasificar el parche en la categoría estimada mayoritaria, la cual toma parches de 48x48x4 y los clasifica en un vector de probabilidades de 1x19. Posteriormente se diseñó un modelo
basado en una \acrshort{cnn} con estructura U-Net, la cual clasifica los parches de 48x48 a nivel de píxel, asignando a cada uno una de las 19 etiquetas propuestas por BigEarthNet. Este trabajo 
pretende acercarse al nivel de BigEarthNet sobre suelo español, con el fin de actualizar en un futuro los usos de suelo a nivel nacional de forma más rápida y eficiente de forma automática
o al menos como un apoyo.
}

\keywordses{Segmentación semántica, \acrfull{cnn}, Sentinel-2, CORINE Land Cover, Usos de suelo, Inteligencia artificial}

\MediaOptionLogicBlank

\pdfbookmark[1]{Abstract}{abstract}
\chapter*{Abstract}
\guideinfo{The objective of this project is the design and implementation of an \acrfull{ai} model capable of automating the semantic segmentation of land cover and 
land use in Spain using Sentinel-2 satellite imagery. The proposed methodological approach involved extracting the satellite images through Google Earth Engine and downloading 
the national level classified land use database provided by \acrfull{corinein} Land Cover Spain. These data were analyzed using the Rasterio library to verify that they shared the 
same \acrfull{crs}. Additionally, the vector information within the geopackage was rasterized and projected onto the EPSG:25830 system to align it with the satellite images at a 
spatial resolution of 10 meters per pixel, and finally divide the macro patches into more manageable 50x50 patches.
\newline
For the development of the neural network, two distinct models were proposed: First, a \acrfull{cnnin} coupled with a \acrfull{mlpin} aimed at classifying each patch into its estimated 
majority category, which takes 48x48x4 patches and maps them into a 1x19 probability vector. Subsequently, a model based on a \acrshort{cnnin} with a U-Net architecture was designed 
to perform pixel-level classification on the 48x48 patches, assigning each coordinate one of the 19 labels proposed by BigEarthNet. This work aims to approach the benchmark performance 
of BigEarthNet on Spanish soil, with the ultimate purpose of enabling faster, more efficient, and automated updates of national land use maps in the future, or at least serving as a support tool.
}

\keywordsen{Semantic segmentation, \acrfull{cnn}, Sentinel-2, CORINE Land Cover, Land use, Artificial intelligence}

\MediaOptionLogicBlank